\section{The rank-six cusp lattice inside the nine-dimensional cone program}
\label{rsl:sec:rank-six}

The rank-three toroidal construction contains a free rank-six lattice before
any finite quotient is taken.  This section identifies that lattice, its
integral symmetric-group action, its monodromy, and the exact part of it seen
by the Hermitian Gram boundary.  It also separates it from the different
rank-six group which classifies bundle defects on the regular locus of the
reconstructed \(S^6\).

\subsection{The two rank-six groups have different types}
\label{rsl:subsec:two-lattices}

Retain the notation of \cref{subsec:teb-rank-three}.  The lattice native to
the rank-three toroidal cusp is
\begin{equation}
 \Lambda_{\rm tor}
 :=H_1(T^6;\mathbb Z)
 =L_3\oplus\overline L_3
 =\mathbb Z\langle e_1,e_2,e_3\rangle
  \oplus\mathbb Z\langle\bar e_1,\bar e_2,\bar e_3\rangle.
 \label{rsl:eq:tor-lattice}
\end{equation}
Its first summand consists of angular cycles and its second summand consists
of deck cycles.  The cusp monodromy and filling map are
\begin{equation}
 M_3(a,b)=(a+b,b),
 \qquad
 c_3(a,b)=b.
 \label{rsl:eq:monodromy-filling}
\end{equation}

The bundle-defect lattice on the original threefold has a different source
and a different cohomological degree.  With
\(D=W\sqcup S_1\sqcup S_2\) and \(U=X\setminus D\), it is
\begin{equation}
 L_{\rm def}:=H^4(U;\mathbb Z)
 \cong H_1(W;\mathbb Z)\oplus H_1(S_1;\mathbb Z)
                    \oplus H_1(S_2;\mathbb Z)
 \cong(\mathbb Z^2)^3.
 \label{rsl:eq:defect-lattice}
\end{equation}
The map \(\Phi=P\oplus Q\) below is defined on
\(\Lambda_{\rm tor}\), whereas the fibre-degree character of
\cref{zdc:thm:fibre-degree} is defined on \(L_{\rm def}\).  Equality of their
ranks supplies no map between them.  In particular, neither group is the
cyclic group \(\mathbb Z/6\mathbb Z\).

There is also no component-permuting symmetric-group action on
\(L_{\rm def}\) supplied by the marked candidate fibration.  The three
deleted fibres have, respectively,
unipotent, order-three, and order-four local monodromy.  An automorphism of
the marked fibration which permutes the punctures must preserve these
conjugacy-order types and hence fixes all three punctures.  One may impose an
abstract permutation of the three displayed summands in
\eqref{rsl:eq:defect-lattice} after choosing three identifications with
\(\mathbb Z^2\), but that is a re-marking, not an automorphism furnished by
the geometry.  This argument does not classify automorphisms fixing the base
pointwise or internal actions on the individual summands; in particular it
does not assert that the abstract group admits no \(S_3\)-action.

\subsection{The integral \texorpdfstring{\(S_3\)}{S3}-module of the cusp}
\label{rsl:subsec:s3-module}

Let \(R_\sigma\) permute the three standard basis vectors by
\(R_\sigma e_i=e_{\sigma(i)}\), and let \(S_3\) act diagonally on
\(\Lambda_{\rm tor}=L_3\oplus\overline L_3\).  Put
\begin{equation}
 k=e_1+e_2+e_3,
 \qquad
 K=\mathbb Z(k,0)\oplus\mathbb Z(0,\bar k),
 \quad \bar k=\bar e_1+\bar e_2+\bar e_3.
 \label{rsl:eq:kernel-k}
\end{equation}
The group \(S_3\) fixes \(K\) pointwise, commutes with \(M_3\), and makes
\(c_3\) equivariant for the permutation action on \(\overline L_3\):
\(c_3(R_\sigma a,R_\sigma b)=R_\sigma c_3(a,b)\).

Recall
\begin{equation}
 P=\begin{pmatrix}1&-1&0\\0&1&-1\end{pmatrix},
 \qquad
 Q=\begin{pmatrix}0&-1&1\\1&-1&0\end{pmatrix},
 \qquad
 \Phi=P\oplus Q.
 \label{rsl:eq:PQ}
\end{equation}
For \(\sigma\in S_3\), define integral actions on \(A_F\) and \(B_F\) by
\begin{equation}
 \rho_A(\sigma)(Px)=P(R_\sigma x),
 \qquad
 \rho_B(\sigma)(Qx)=Q(R_\sigma x).
 \label{rsl:eq:target-actions}
\end{equation}
These definitions are well posed because
\(\ker P=\ker Q=\mathbb Zk\).

\begin{theorem}[The exact rank-six \(S_3\)-module sequence]
\label{rsl:thm:s3-sequence}
The map \(\Phi\) fits into an exact sequence of integral
\(S_3\)-modules
\begin{equation}
 0\longrightarrow K\longrightarrow\Lambda_{\rm tor}
 \xrightarrow{\ \Phi\ }A_F\oplus B_F\longrightarrow0.
 \label{rsl:eq:s3-sequence}
\end{equation}
It is simultaneously equivariant for the cusp monodromies:
\begin{equation}
 \Phi M_3=M_F\Phi,
 \qquad
 M_3|_K=\begin{pmatrix}1&1\\0&1\end{pmatrix}
 \quad\text{in the basis }((k,0),(0,\bar k)).
 \label{rsl:eq:kernel-monodromy}
\end{equation}
As a sequence of free abelian groups with monodromy it splits integrally, as
shown by \cref{prop:teb-explicit-conjugacy}.  As a sequence of integral
\(S_3\)-modules it does not split.  After tensoring with \(\mathbb Q\),
\begin{equation}
 \Lambda_{\rm tor,\mathbb Q}\cong
 2\mathbf 1\oplus2V_{\rm std},
 \qquad
 (A_F\oplus B_F)_{\mathbb Q}\cong2V_{\rm std},
 \qquad K_{\mathbb Q}\cong2\mathbf1.
 \label{rsl:eq:rational-decomposition}
\end{equation}
\end{theorem}

\begin{proof}
The kernels of \(P\) and \(Q\) are both the primitive diagonal lattice
\(\mathbb Zk\), and both maps are surjective.  This proves exactness.
Equations \eqref{rsl:eq:target-actions} make \(P\), \(Q\), and hence
\(\Phi\), equivariant.  Since \(R_\sigma(a+b)=R_\sigma a+R_\sigma b\), the
source action commutes with \(M_3\).  Moreover
\[
 M_3(xk,y\bar k)=((x+y)k,y\bar k),
\]
which gives \eqref{rsl:eq:kernel-monodromy}.  The target equivariance with
\(M_F\) is equivalently the already verified identity
\(\Phi M_3=M_F\Phi\).

It remains to prove the integral non-splitting.  Any \(S_3\)-equivariant
homomorphism \(r:\mathbb Z^3\to\mathbb Z\), where the target is trivial, has
the same value \(a\in\mathbb Z\) on all three standard basis vectors.  Thus
\(r(k)=3a\), so no such map retracts the inclusion
\(\mathbb Zk\hookrightarrow\mathbb Z^3\).  For two copies, every entry of
the matrix obtained by restricting an equivariant map
\((\mathbb Z^3)^2\to\mathbb Z^2\) to \((\mathbb Zk)^2\) is divisible by
three.  That restriction cannot be the identity matrix.  Hence
\eqref{rsl:eq:s3-sequence} does not split over \(\mathbb Z[S_3]\).

Over \(\mathbb Q\), the invariant diagonal line has the equivariant
complement \(\{x_1+x_2+x_3=0\}\), the standard representation.  Applying
this to both copies gives \eqref{rsl:eq:rational-decomposition}.
\end{proof}

For reference, if \(s_{12}\) and \(s_{23}\) are the adjacent
transpositions, the induced target matrices are
\begin{equation}
 \rho_A(s_{12})=\begin{pmatrix}-1&0\\1&1\end{pmatrix},
 \quad
 \rho_A(s_{23})=\begin{pmatrix}1&1\\0&-1\end{pmatrix},
 \label{rsl:eq:rho-A}
\end{equation}
and
\begin{equation}
 \rho_B(s_{12})=\begin{pmatrix}1&-1\\0&-1\end{pmatrix},
 \quad
 \rho_B(s_{23})=\begin{pmatrix}-1&0\\-1&1\end{pmatrix}.
 \label{rsl:eq:rho-B}
\end{equation}
Direct multiplication gives
\begin{equation}
 \rho_A(\sigma)P=PR_\sigma,
 \qquad
 \rho_B(\sigma)Q=QR_\sigma,
 \qquad
 \rho_A(\sigma)B_0=B_0\rho_B(\sigma).
 \label{rsl:eq:rho-intertwiners}
\end{equation}

The quotient \(\mathbb Z^3/\mathbb Zk\) is the \(A_2\) weight lattice, not
the \(A_2\) root lattice.  Indeed, the augmentation-zero lattice is generated
by \(e_1-e_2\) and \(e_2-e_3\), whose images under \(P\) are the columns of
\begin{equation}
 \begin{pmatrix}2&-1\\-1&2\end{pmatrix}.
 \label{rsl:eq:root-index}
\end{equation}
Its determinant is three.  Thus the root lattice is the index-three
sublattice of \(A_F\) generated by the differences of the three weights
\(n_1,n_2,n_3\).  The integral non-splitting in
\cref{rsl:thm:s3-sequence} is the same index-three phenomenon viewed before
passing to the quotient.

\subsection{The six homology generators are not the six chambers}
\label{rsl:subsec:three-sixes}

There are three related but distinct appearances of the integer six.

\begin{theorem}[The three sixes]
\label{rsl:thm:three-sixes}
For the rank-three toroidal cusp and every injective integral shear
\(\mathcal B:\overline L_3\hookrightarrow L_3\),
\begin{align}
 \operatorname{rk}\Lambda_{\rm tor}&=2\operatorname{rk}L_3=6,
 \label{rsl:eq:six-rank}\\
 \#\{\text{Kuhn maximal cells modulo }L_3\}&=|S_3|=3!=6,
 \label{rsl:eq:six-cells}\\
 \chi(W_{3,\mathcal B})
 &=|S_3|[L_3:\mathcal B\overline L_3]
   =3!\,|\det\mathcal B|.
 \label{rsl:eq:six-euler}
\end{align}
The equalities \eqref{rsl:eq:six-cells} and \eqref{rsl:eq:six-euler} have the
same toric source: the maximal cells are indexed by orderings of three angular
coordinates.  Equation \eqref{rsl:eq:six-rank} has a different source: first
homology contains one angular and one deck copy of the rank-three lattice.
For a general shear \(\mathcal B\), this cell count does not assert that the
diagonal \(S_3\)-action descends to the sheared quotient; that requires the
additional commutation equations
\(R_\sigma\mathcal B=\mathcal B R_\sigma\).

In particular, the six standard basis vectors of \(\Lambda_{\rm tor}\) are
not \(S_3\)-equivariantly identifiable with the six Kuhn cells.  The basis
vectors form two orbits of size three; the maximal cells form one free orbit
of size six.
\end{theorem}

\begin{proof}
The rank formula follows from \eqref{rsl:eq:tor-lattice}.  A Kuhn cell in one
lattice cube is labelled by the order in which the three coordinate
increments are taken, hence by a unique element of \(S_3\).  Left
multiplication is free and transitive on these labels, proving
\eqref{rsl:eq:six-cells}.  The normalization and orbit calculation in
\cref{hrftd:cor:euler-shear-determinant} gives
\eqref{rsl:eq:six-euler}.

Under the diagonal permutation action, the angular basis
\(\{e_1,e_2,e_3\}\) and the deck basis
\(\{\bar e_1,\bar e_2,\bar e_3\}\) are two distinct transitive three-element
sets.  The cell labels form the regular six-element \(S_3\)-set.  No
equivariant bijection exists.  Equivalently, over \(\mathbb Q\) the character
of the basis permutation module on the conjugacy classes
\((1),(12),(123)\) is
\[
 (6,2,0),
\]
whereas the regular representation has character \((6,0,0)\).  The
representations are not isomorphic even over \(\mathbb Q\).
\end{proof}

This distinction also locates the chamber construction: the Kuhn fan uses the
angular summand \(L_3\), while the deck summand is needed for periods,
monodromy, and filling.  The map \(P\) produces the chamber fan; the different
map \(Q\) makes the filling square commute.  Replacing both by a single
untyped ``six-direction map'' loses precisely the information that makes the
cusp morphism valid.

\subsection{The exact map into the Gram boundary}
\label{rsl:subsec:gram-map}

Let
\[
 \mathfrak a
 =\{\operatorname{diag}(\lambda_1,\lambda_2,\lambda_3):
       \lambda_1+\lambda_2+\lambda_3=0\}
 \subset\operatorname{Herm}_3^0(\mathbb C),
\]
let \(\jmath:A_{F,\mathbb R}\to\mathfrak a\) be the marked isomorphism
\(\jmath(n_i)=E_{ii}-I_3/3\), and let
\[
 \Phi_{\rm G}(H)=H-\lambda_{\min}(H)I_3.
\]
Write \(\operatorname{pr}_{\rm ang}(a,b)=a\).  The maps which actually exist
are
\begin{equation}
 \Psi=\jmath\circ P\circ\operatorname{pr}_{\rm ang}:
 \Lambda_{\rm tor,\mathbb R}\longrightarrow\mathfrak a,
 \label{rsl:eq:linear-gram-map}
\end{equation}
and
\begin{equation}
 \Gamma=\Phi_{\rm G}\circ\Psi:
 \Lambda_{\rm tor,\mathbb R}\longrightarrow
 \{H\in\operatorname{Herm}_3(\mathbb C):H\succeq0,\ \det H=0\}.
 \label{rsl:eq:piecewise-gram-map}
\end{equation}

\begin{proposition}[Exact scope of the Gram-boundary map]
\label{rsl:prop:gram-scope}
The map \(\Psi\) is linear and \(S_3\)-equivariant, with
\begin{equation}
 \ker\Psi=\mathbb R(k,0)\oplus(0,\overline L_{3,\mathbb R}).
 \label{rsl:eq:psi-kernel}
\end{equation}
The map \(\Gamma\) is continuous, positively homogeneous,
\(S_3\)-equivariant, and piecewise linear.  On the six projected Kuhn cones it
is exactly the six diagonal Weyl-chamber map of
\cref{prop:teb-weyl-chambers}.

The map \(\Gamma\) is not additive, is not proper, ignores the complete deck
summand, and does not descend to a map from a smooth torus fibre.  Thus it is
an exact continuous piecewise-linear chamber map on a real cocharacter space
into the diagonal part of the nine-dimensional Gram carrier's
eight-dimensional boundary, not a toric fan morphism and not a morphism from
the toroidal total space to the Hermitian Gram cone.
\end{proposition}

\begin{proof}
Linearity and equivariance of \(\Psi\) follow from
\eqref{rsl:eq:rho-intertwiners} and the Weyl permutation of diagonal entries.
Since \(\jmath\) is injective and \(\ker P=\mathbb Rk\), its kernel is exactly
\eqref{rsl:eq:psi-kernel}.  The least-eigenvalue function is continuous,
positively homogeneous, and linear on each closed Weyl chamber, which proves
the corresponding assertions for \(\Gamma\).  The chamber formula is
\cref{prop:teb-weyl-chambers}.

The function \(H\mapsto H-\lambda_{\min}(H)I_3\) is not additive across
different Weyl chambers.  The kernel in \eqref{rsl:eq:psi-kernel} has positive
dimension, so \(\Gamma\) is not proper.  It ignores the deck summand by
definition.  Finally, translation by the angular lattice does not leave it
invariant: for example \(\Gamma(e_1,0)\neq\Gamma(0,0)\).  Hence it cannot
descend through the quotient of the universal cover by
\(\Lambda_{\rm tor}\).
\end{proof}

There is one further integral check on what the elliptic quotient does not
preserve.  On \(\Lambda_{\rm tor}\), define the standard alternating form
\begin{equation}
 E_3((a,b),(a',b'))=a^{\mathsf T}b'-b^{\mathsf T}a'.
 \label{rsl:eq:standard-alternating}
\end{equation}
It is invariant under \(S_3\) and \(M_3\), but on the kernel it satisfies
\begin{equation}
 E_3((k,0),(0,\bar k))=3.
 \label{rsl:eq:kernel-degree-three}
\end{equation}
Every alternating form pulled back along \(\Phi\) vanishes whenever one
argument lies in \(K\).  Therefore \(E_3\) is not the pullback of an
alternating form on \(A_F\oplus B_F\).  In particular, the holomorphic
elliptic quotient in \cref{thm:teb-holomorphic-cusp-bridge} is not a morphism
preserving this chosen principal lattice form.  This is an integral statement
about \eqref{rsl:eq:standard-alternating}; it does not assert that either
nonalgebraic torus family carries a polarization.

\begin{auditrecord}[Rank-six conclusion]
The native higher-rung object is the monodromy lattice
\(\Lambda_{\rm tor}=H_1(T^6;\mathbb Z)\), not a cyclic group.  Its exact
relation to the nine-dimensional Hermitian construction factors through the
angular quotient to the \(A_2\) weight lattice and then through the marked
Cartan/Weyl fan.  The deck half is indispensable for the period and filling
squares but invisible to the Gram-boundary chamber map.  The separate
six-coordinate group \(H^4(U;\mathbb Z)\) remains the correct classifier for
bundle defects on the original \(S^6\) regular locus; no identification of
these two rank-six groups is used.
\end{auditrecord}
